\section{Primitive coordinates at one complete output key}
\label{sec:primitive-coordinates}

For every ordered pair \(m\ne n\), define
\begin{equation}
 g_{mn}=(d_m,d_n),\qquad
 a_{mn}=\frac{d_m}{g_{mn}},\qquad
 b_{mn}=\frac{d_n}{g_{mn}}.
 \label{eq:row-divisor-split}
\end{equation}
The three factors \(g_{mn},a_{mn},b_{mn}\) are squarefree and pairwise
coprime, and reversal gives
\[
 g_{nm}=g_{mn},\qquad a_{nm}=b_{mn},\qquad b_{nm}=a_{mn}.
\]
Fix distinct active rows \(m<n\) and abbreviate
\[
 g=g_{mn},\qquad a=a_{mn},\qquad b=b_{mn}.
\]
Put
\[
 \Delta_{mn}=h_0(n-m)\ne0.
\]

\begin{definition}[Direct-entry cofactor coordinates]
\label{def:direct-entry-coordinates}
For \(i\in\cI_{mn}\), define
\begin{equation}
 q_i=(S_{i,m},S_{i,n}),
 \qquad
 S_{i,m}=q_i u_i,
 \qquad
 S_{i,n}=q_i v_i.
 \label{eq:direct-entry-coordinates}
\end{equation}
\end{definition}

\begin{theorem}[Canonical primitive key coordinates]
\label{thm:canonical-primitive-key}
At every actual key \(i\in\cI_{mn}\):
\begin{enumerate}[label=\textup{(\roman*)},leftmargin=3em]
 \item \(q_i,u_i,v_i\) are squarefree and pairwise coprime;
 \item \(q_i\) is coprime to \(gab\), \(u_i\) is coprime to \(ga\), and
       \(v_i\) is coprime to \(gb\);
 \item
       \begin{equation}
       \boxed{q_i\mid |\Delta_{mn}|;}
       \label{eq:direct-q-divides-Delta}
       \end{equation}
 \item the two largest-prime labels satisfy
       \begin{equation}
       \boxed{
       n u_i Q_{i,m}-m v_iQ_{i,n}=\frac{\Delta_{mn}}{q_i};}
       \label{eq:direct-prime-line}
       \end{equation}
 \item the residual labels are equal exactly in one primitive direction:
       \begin{equation}
       \boxed{
       s_{i,m}=s_{i,n}
       \quad\Longleftrightarrow\quad
       (u_i,v_i)=(b,a).}
       \label{eq:unique-residual-direction}
       \end{equation}
\end{enumerate}
\end{theorem}

\begin{proof}
The residual labels are
\[
 s_{i,m}=gaq_i u_i,\qquad
 s_{i,n}=gbq_i v_i.
\]
Both are squarefree by \eqref{eq:squarefree-key-carrier}.  Together with
\((u_i,v_i)=1\), which follows from the definition of \(q_i\), this gives
the stated squarefreeness and coprimality: squarefreeness of
\(gaq_iu_i\) makes \(g,a,q_i,u_i\) pairwise coprime, and squarefreeness
of \(gbq_iv_i\) does the same for \(g,b,q_i,v_i\).

The common time in the complete key gives
\begin{align*}
 nN_{i,m}-mN_{i,n}
 &=
 n(mj(i)+h_0)-m(nj(i)+h_0)\\
 &=h_0(n-m)=\Delta_{mn}.
\end{align*}
Substitute \(N_{i,m}=Q_{i,m}q_i u_i\) and
\(N_{i,n}=Q_{i,n}q_i v_i\).  This yields
\[
 q_i\bigl(nu_iQ_{i,m}-mv_iQ_{i,n}\bigr)=\Delta_{mn},
\]
proving \eqref{eq:direct-q-divides-Delta} and
\eqref{eq:direct-prime-line}.

Finally,
\[
 s_{i,m}=s_{i,n}
 \quad\Longleftrightarrow\quad
 au_i=bv_i.
\]
If \(au_i=bv_i\), then \((a,b)=1\) gives
\(u_i=bk\), \(v_i=a\ell\).  The equality gives \(k=\ell\), while
\((u_i,v_i)=1\) forces \(k=\ell=1\).  The converse is immediate.
\end{proof}

\begin{remark}[No anchor is used]
\label{rem:no-anchor-needed}
The coordinate \(q_i\) is recovered from the two cofactors at the same
complete output key.  A same-residual return key is not part of
\cref{def:direct-entry-coordinates}.  The row-divisor split only marks
which point of the full primitive plane belongs to \(R\).
\end{remark}

\begin{corollary}[Divisor-many content slices]
\label{cor:divisor-many-direct-slices}
For fixed \(m,n,h_0\), every actual entry key belongs to one of at most
\(\tau(|\Delta_{mn}|)\) content slices.  This divisor bound controls
multiplicity only; it gives no cancellation among literal slice
amplitudes.
\end{corollary}

\begin{remark}[Actual support is smaller]
The arithmetic conditions above are necessary, not sufficient.  Primality
of \(Q_{i,m},Q_{i,n}\), their largest-prime status, the physical interval,
complete-key incidence, sides, masks, profiles, and endpoints may delete
points from a primitive direction.  They are retained inside the literal
kernel defined below.
\end{remark}
